\documentclass[a4paper]{article}
\usepackage{amsfonts}
\usepackage{amsmath}
\usepackage{amssymb}
\usepackage{graphicx}     

\newcommand{\dint}{\displaystyle\int}
\newcommand{\llim}{\lim\limits}

\begin{document}
  
\noindent \large Auteur: Ayoub Jaaouan          \\
\noindent \large Studierichting: Informatica   \\ 
\noindent \large Oefeningenreeks 1C nr. 2.7   \\

\medskip

\normalsize

$\left(-3, \sqrt{3}\right)$

$\left \{\begin{array}{l} r = \sqrt{(-3)^2+ (\sqrt{3})^2} \\ \tan \theta = \dfrac{-\sqrt{3}}{3} \end{array}\right. 
\Rightarrow \left \{\begin{array}{l} r = \sqrt{9+ 3} \\ \tan \theta = \dfrac{-\sqrt{3}}{3} \end{array}\right.$ \\



$\Rightarrow r = \sqrt{(-3)^2+ (\sqrt{3})^2} = \sqrt{9+ 3} = \sqrt{12}$\\

$\Rightarrow \tan \theta = \dfrac{-\sqrt{3}}{3} \Rightarrow \operatorname{Bgtan}\left(\dfrac{-\sqrt{3}}{3}\right) = \dfrac{5\pi}{6}  $\\



\textit{\textbf{Resultaat:}} $\left(\sqrt{12}, \dfrac{5\pi}{6}\right)$


\begin{thebibliography}{999}
%\bibitem{a} Referentie
\end{thebibliography}
\end{document}